\section{From the Type-I level forms to weighted segment bounds}
\label{sec:segment-weighted}

\status{The actual integral level comparison, block-lattice conversion,
terminal Gram window and weighted summation are proved here.
The final digit-class estimate uses TP's proved multiple-tie mean and
explicitly retains the coarse rank bound (SW8.1) and the stated numerical
conditions on the parameter box. No global dimension-density theorem is
inferred from the summation alone.}

\paragraph{SW1: levels are a direct sum, and their integral span is saturated.}
Use TI1's variables, and write
$j=(j_1,\ldots,j_t,j_+)$ for a level.  Before multiplying by
$\prod E_i^{j_i}E_+^{j_+}\prod Z_i^{z_i}Z_+^{z_+}$, its lattice is exactly
\[
 L_j=Y_{\mathbb Z_p}(d(j))\otimes
           \mathbb Z_p[\xi]_{B-j_+},\qquad
 \operatorname{rank}L_j=\dim Y(d(j))\binom{B-j_++r-1}{r-1}.
 \tag{SW1.1}
\]
The tensor product is embedded by multiplication in disjoint variables.
Its integral description follows coefficient by coefficient in $\xi$.
For a finite relation between different levels, take the lexicographically
largest $S$-degree $j$. Its coefficient is the nonzero monomial
$\prod Z_i^{z_i+j_i}Z_+^{z_++j_+}$ times that level's coefficient polynomial;
no smaller level contributes. Thus the relation is zero level by level.
The same argument proves saturation: if the total polynomial is integral,
extracting this coefficient makes its chain coefficients integral and hence
members of $Y_{\mathbb Z_p}$, subtract the corresponding integral level,
and descend. Formation of $Y$ commutes with field extension: in its finite
multidegree space, each prefix condition is the common kernel of its
forbidden-weight coefficient maps after words in that prefix's own
generators. Each decreasing kernel sequence stabilizes; intersect the
finitely many resulting kernels. Thus the actual block
rank, also over the ramified coefficient field, equals the sum in (SW1.1).
This argument uses the actual $S$ and $Z$ variables, not a dimension count
for a substituted ambient space.

\paragraph{SW2: the integral harmonic change of levels.}
For one pair of residual degrees $(o,n)$ and shift $z$, put
$\widetilde E=ZS+\sum_c w_ce_c$ and
$\widetilde F=\partial_Z\partial_S+\sum_c\partial_{w_c}\partial_{e_c}$.
For a polynomial $f$ independent of $Z,S$, define
\[
 P_{o,n,z}f=\sum_{\ell=0}^{\min(o,n)}
 \frac{(-1)^\ell\widetilde E^\ell Z^zF^\ell f}
 {\ell!^2\binom{o+n+r+z-1}{\ell}}.
 \tag{SW2.1}
\]
The commutator $[\widetilde F,\widetilde E]
=r+1+\mathrm{Euler}$ gives
$\widetilde F P_{o,n,z}f=0$. Its inverse expansion is
\[
 Z^zf=\sum_{\ell=0}^{\min(o,n)}
 \widetilde E^\ell P_{o-\ell,n-\ell,z}
 \left(\frac{F^\ell f}
 {\ell!^2\binom{o+n+r+z-\ell}{\ell}}\right).
 \tag{SW2.2}
\]
These finite identities follow by collecting $\widetilde E^\ell F^\ell$
and the falling-factorial binomial identity; the companion gives the
coefficient cancellation explicitly.

Both expansions cost at most
$u=r\lfloor\log_p(2A)\rfloor+\lfloor\log_p(6A+r)\rfloor$
when the relevant row degrees and shifts are at most $2A$.
Indeed $p^{r\lfloor\log_p(2A)\rfloor}F^\ell/\ell!^2$ preserves the
appropriate integral residual lattice, by multinomial carries, and the
remaining inverse binomial costs at most the second logarithm.
The prefix conditions are preserved: a vector-row derivative and its
coordinate companions form a stable space whose first weight does not
increase; a complete vector--covector contraction commutes with the prefix
action and preserves weight. If the prefix cuts the pair, only its vector
row is inside, so the first observation applies.
This proves the source's \texttt{HypJoin} and \texttt{HypU}.
Different pairs have disjoint variables. Applying (SW2.1)--(SW2.2) successively
therefore proves, for the actual original and harmonic level lattices
$SS_{\mathbb Z_p},PS_{\mathbb Z_p}$,
\[
 p^U PS_{\mathbb Z_p}\subset SS_{\mathbb Z_p},\qquad
 p^U SS_{\mathbb Z_p}\subset PS_{\mathbb Z_p},\qquad
 U=(t+1)\left\lceil(r+1)(\log_p(8A+r)+1)\right\rceil.
 \tag{SW2.3}
\]
Their rational spans agree. Since $\widetilde E$ and $\widetilde F$ are
Fock adjoints, different harmonic levels are orthogonal. On level $j$,
the unnormalized harmonic map $\Theta_j$ has pairing
\[
 B(\Theta_jf,\Theta_jg)=c_j\,B(\Phi_{\rm all}f,g),\quad
 c_j=\prod_{\nu=1}^{t+1}
 j_\nu!^2 z_\nu!\binom{o_\nu+n_\nu+r+z_\nu+j_\nu}{j_\nu}.
 \tag{SW2.4}
\]
Here $o_\nu,n_\nu$ are the residual degrees at that pair.
The norm identity follows from
$\widetilde F\widetilde E^\ell h
=\ell(d_h+r+\ell)\widetilde E^{\ell-1}h$
for harmonic $h$, and
$B(Pf,Pg)=z!B(f,\Phi_{o,n,z}g)$.

\paragraph{SW3: convert the good level factors into the actual block lattice.}
Let $d_j!=\prod_a d_a(j)!$ include the terminal $\xi$ degree, and let
$b_j(D)$ be the bad Smith count of
$B(\Phi_{\rm all}\,\cdot,\cdot)/d_j!$ on $L_j$ at depth $D\in\mathbb Z$.
Choose dual families for all good factors, move them by $p^U\Theta_j$,
and multiply the partner on level $j$ by
$p^M/(c_jd_j!)$, where $M=\max_j v_p(c_jd_j!)$.
The multiplier is integral; (SW2.3)--(SW2.4) give one common pairing
$p^e\delta_{ab}$, with $e=D+2U+M$, and at least
$\operatorname{rank}SS_{\mathbb Z_p}-\sum_jb_j(D)$ pairs.

Let $\Lambda=\log_p(8a+r)+1$, $A=(3s+2)a$, $t\le s$, $r\ge4$ and
$p\ge(r+1)^2$, with $2s+2\le r$. Legendre's formula, applied to the
actual row multidegrees, gives for every supported monomial $\mu$ and level $j$
\[
 v_p(c_jd_j!)\le v_p(\mu!)+
       \big((N+1)r+3(t+1)\big)\lfloor\log_p(8A+r)\rfloor,
 \qquad N=2s+1.
 \tag{SW3.1}
\]
All levels have the same extended row multidegrees: each $Z$ belongs to
its vector row and each $S$ to its covector row. There are
$(N+1)r+2(t+1)$ variables; bounding multinomial carries in these rows
gives that many logarithms. The remaining $t+1$ logarithms bound the
binomials in $c_j$. The factorial part is bounded by the row factorial
product because $j!^2z!o!n!$ divides $(o+z+j)!(n+j)!$.
Since $\log_p(8A+r)\le\Lambda$, it follows that
\[
 e\le v_p(\mu!)+D+\ell_{\rm lev},\qquad
 \ell_{\rm lev}=10(s+1)(r+1)\Lambda+\kappa_0,\quad \kappa_0\ge0.
 \tag{SW3.2}
\]

Here is the exact ramified-lattice conversion. Write $s_b$ for TI1's
factorial constant, $w^2=p$, $F_0=\lfloor D+\ell_{\rm lev}\rfloor$, and
$q=F_0-e-v_p(s_b)$. Multiply both families by $w^q$.
For each coefficient, (SW3.2) gives
$q+v_p(\mu!)+v_p(s_b)\ge0$, precisely TI1's weighted integrality condition.
Their block pairing has valuation
$v_p(s_b)+q+e=F_0$. Saturating the first-family span only improves the
primitive-test bound: a primitive vector of its saturation needs at
worst division by a power of $w$, which lowers the valuation of a suitable
pairing. Thus the source's actual \texttt{BlockBoundUniformiser} holds,
with rank loss $\sum_jb_j(D)$ and depth $D+\ell_{\rm lev}$.

\paragraph{SW4: tensoring and diagonalizing the intermediate form.}
Before the terminal contraction the form is
$b_0=B(\Phi_{\rm ties}\,\cdot,\cdot)/d_j!$ on (SW1.1).
If $B(\Phi_{\rm ties}x_i,y_k)/\prod_a d_a!=p^D\delta_{ik}$,
then
$x_i\xi^\alpha$ and
$\big((B-j_+)!/\alpha!\big)y_i\xi^\alpha$
are integral dual families for $b_0$.
Their pairing is $p^D\delta_{ik}\delta_{\alpha\beta}$.
Consequently
\[
 \operatorname{bad}(b_0,L_j,D)
 \le\operatorname{bad}(b_Y,Y_{\mathbb Z_p},D)
          \binom{B-j_++r-1}{r-1}.
 \tag{SW4.1}
\]
The disjoint tie operators are self-adjoint and commute, so $b_0$ is
symmetric. Over $\mathbb Z_p$ with $p\ne2$, every finite free lattice
with symmetric rational form has an orthogonal integral basis.
Choose a vector with least nonzero pairing valuation. If a minimal
entry is off the diagonal, the sum of its two basis vectors has that
valuation since $2$ is a unit. Pairing divided by its diagonal entry is
integral, giving an orthogonal direct summand; iterate, then retain the
radical. The good basis vectors are exactly those whose diagonal
valuations are at most $D$. This justifies the orthogonal-family step
without assuming a self-dual lattice.

\paragraph{SW5: the terminal Gram window is another polynomial average.}
Fix an orthogonal good family $f_1,\ldots,f_n$ for $b_0$, with nonzero
diagonal entries $d_i$ and $v_p(d_i)\le D$. Let $o,n_\xi$ be the remaining
degrees, $J=\min(o,n_\xi)$, $q=o+n_\xi+r-1$, and vary the terminal shift $z$.
For
$B_z(f,g)=B(\Phi_{o,n_\xi,z}\Phi_{\rm ties}f,g)/d_j!$, set
\[
 \begin{gathered}
 M(z)=p^{x_0}\binom{q+z}{J}\big(B_z(f_i,f_k)\big)_{i,k},\\
 x_0=\big((t+1)(r+1)+(N+1)r+2(t+1)\big)L,\quad
 L=\lfloor\log_p(8A+r)\rfloor+1.
 \end{gathered}
 \tag{SW5.1}
\]
It is an integral binomial-polynomial matrix of degree at most $J$.
Its leading matrix is $p^{x_0}\operatorname{diag}(d_i)/J!$.
To see this, multiply the projector formula by $\binom{q+z}J$:
its $\ell$th term becomes
\[
 (-1)^\ell\binom{q+z-\ell}{J-\ell}
 \frac{B(E^\ell F^\ell\Phi_{\rm ties}f_i,f_k)}
      {\ell!^2\binom J\ell\,d_j!}.
 \tag{SW5.2}
\]
The first $t+1$ operations cost $(t+1)(r+1)L$ powers of $p$;
normalized Fock pairing costs at most $((N+1)r+2(t+1))L$ by row
multinomial carries. Vandermonde's identity gives integral binomial
coefficients in $z$. Only $\ell=0$ has degree $J$, proving the leading formula.
Its determinant has degree $Jn$, and Legendre gives
$v_p(\operatorname{lc}\det M)+Jn/(p-1)\le n(x_0+D+L)$.
The scalar in (SW5.1) is nonzero and integral for $z\ge0$; multiplying
by it can only increase bad Smith counts. TI5 therefore yields
\[
 \frac1W\sum_{i=0}^{W-1}\operatorname{bad}(B_{z_0+i},\langle f\rangle,k)
 \le n\left\{\frac JW+
       \frac{x_0+D+L+Jh/W}{k+1}\right\},
 \quad W,J+1\le p^h,\quad k\ge0.
 \tag{SW5.3}
\]

In the actual window put
$\kappa=3(s+2)r^2\Lambda$,
$D_1=t(Y\kappa+C_M\kappa)+C_M\kappa$, $C_M\ge1$,
$D=\lfloor D_1\rfloor$, $k=\lfloor YD_1\rfloor$.
The literal numerical inputs are
$W\ge a/((r+1)Y)$, $W\le a+1$,
$J\le5b'/2$, $pb'<a$, $p\ge(r+1)^2Y^2$, $Y\ge1$, and $o_+\le a$.
They give $J/W\le1/(2Y)$.
Choose $h=\lceil\log_p(a+1)\rceil$; then $h\le\kappa$ and
$x_0+L\le\kappa\le D_1$. Since $k+1>YD_1$, (SW5.3) is at most
$3n/Y$, hence at most the source's $5n/Y$.
Add the discarded orthogonal directions and use (SW4.1). This proves
the actual \texttt{RestPlus} terminal average.

\paragraph{SW6: fixed digit classes have a fixed index set.}
In a compatible digit class, the coarse degree at site $u$ satisfies
$P_u C_u=a_u-\sum_c\mathrm{digit}_{u,c}$.
The literal segment inequalities give
$(r+1)P_u\le b'$, $\sum_c\mathrm{digit}_{u,c}<b'$,
$b_u\le5b'/2$, $a>25b'$, and $a_u\ge a$.
Thus $C_u\ge A_u$, where $A_u$ is the other coarse degree.
Every kept tie order satisfies $j_i\le o_i\le A_u$,
so its extra condition $j_i\le C_u$ is automatic throughout the box.
The early orders also remain legal: their bound
$\mathrm{sg}_u\le2\lambda A_0/(r+1)$, together with
$4\lambda<r+1$, $A_0\le a/Q$ and $P_u\le Q$,
gives $P_u\mathrm{sg}_u\le a/2<P_uC_u$.
The other kept predicates are independent of $a_u$.
The actual set of kept blocks and legal levels is therefore constant.

Let $u_i(a)$ replace the $i$th tied degree by its largest compatible
endpoint. It stays in the same box. Only one factor in the ambient
dimension $m(a)$ changes, from $\binom{x+r}r$ to $\binom{y+r}r$.
Here $x\ge1$, $y\ge x$, and the source bounds
$xP\ge3a/4$, $yP-xP\le a/((r+1)Y)+x_a$,
$rx_a\le a/8$ give $r(y-x)\le3x$. Hence
\[
 \frac{m(u_i(a))}{m(a)}
 =\prod_{h=1}^r\frac{y+h}{x+h}
 \le(y/x)^r\le\exp\!\big(r(y-x)/x\big)\le e^3.
 \tag{SW6.1}
\]

\paragraph{SW7: the numerical loss in the multiple-tie average.}
TP's exact mean theorem gives, for tie $i$ with progression length $Y_i$,
\[
 \epsilon_i=\frac{2S_i}{Y_i}
       +\frac{2\kappa+2S_i\Lambda/Y_i}{Y\kappa}.
 \tag{SW7.1}
\]
The input is the prefix total $S_i$ with the other tied coordinates at
their upper endpoints, not at the current point.
For the actual box the source reduces the check to
\[
 S_iP_i\le\tfrac52b'+(2s+1)\frac{2a}{p},\quad
 Y_iP_i\ge\frac{a}{2(r+1)Y},\quad
 (r+1)Y^2b'\le a/2,\quad (r+1)^2Y^2\le p.
 \tag{SW7.2}
\]
Thus $S_i/Y_i\le7/Y$. Since $24\Lambda\le\kappa$,
$\epsilon_i\le17/Y\le(30+C_M)/Y$.
Compatible residues form progressions of step one in coarse degree:
division by $P_i$ removes the original spacing $P_i$.
Their initial point obeys
$n_{0,i}P_i\ge a-(r+1)P_i-5b'/2$.
Together with the first bound in (SW7.2), $pb'<a$,
$(r+1)P_i\le b'$ and $p\ge4s+8$, this gives $S_i\le n_{0,i}$.
These checks explain the source's \texttt{RestTMSize}/\texttt{RestTMEps}
interfaces; the other degree and logarithm bounds of TP remain part of
the explicitly stated admissible-input package.

\paragraph{SW8: retain the dimension weights when summing.}
Let $A$ be the compatible top-degree box, $W$ the terminal window,
$I$ the fixed set of block--level indices, and $d(z)$ the terminal
monomial dimension. Write $y(a,z)=\dim Y(a,z)$ and $q(a,z)$ for its
bad count. The required coarse rank input is the literal inequality
\[
 \sum_{z\in I}y(a,z)d(z)\le m(a).
 \tag{SW8.1}
\]
SW1 identifies its left side with the sum of the actual block ranks.
The right-side estimate is the upper-rank clause of
\texttt{flipOneThree.adm11}; it is retained here as a separate input.
Suppose, as proved above and in TP under their matching inputs,
\[
 \begin{aligned}
 \sum_{w\in W}b(a,w,z)&\le |W|
       \{q(a,z)d(z)+5y(a,z)d(z)/Y\},\\
 \sum_{a\in A}q(a,z)&\le\frac{30+C_M}{Y}
       \sum_{i=1}^t\sum_{a\in A}y(u_i(a),z).
 \end{aligned}
 \tag{SW8.2}
\]
Sum first in $z$, use (SW8.1) at $u_i(a)$, then (SW6.1):
\[
 \boxed{\quad
 \sum_{a,w,z}b(a,w,z)
 \le \frac{(30+C_M)t e^3+5}{Y}
                |W|\sum_{a\in A}m(a).\quad}
 \tag{SW8.3}
\]
No injectivity of the endpoint replacement is needed: each updated mass
is bounded separately by $e^3m(a)$. This is why the varying multiplicity
weights cause only a constant loss.

\paragraph{SW9: depth arithmetic.}
For $N_1\ge1$, $s\le N_1$ and $Y\ge1$,
$\kappa\le9N_1(r+1)^2\Lambda$.
It follows directly that
\[
 YD_1+\ell_{\rm lev}
 \le C Y^2N_1^2(r+1)^2\Lambda+\kappa_0
 \quad\text{if }C\ge18(1+C_M)+20.
 \tag{SW9.1}
\]
Taking also
$C\ge2((30+C_M)e^3+5)e^3$
matches the source's \texttt{SegmentWeightedCountingBound}.
With TP's literal constant $C_M=1$, these are fixed constants, independent
of the field, digits and polynomial degrees.

\paragraph{SW10: the coarse-degree gap used in the raw lifts.}
Here $k=r+1$, $\lambda=\mathrm{La}$, $\epsilon=N_1\lambda/k$,
$\delta=\mathrm{de}$, and $\beta=\mathrm{be}$ have their literal
\texttt{Std} meanings. Put $Q_j=Pp^d$, $A_0=\lfloor(a-E_0)/Q_j\rfloor$,
and $A=(g-\sum\text{digits})/P$. Then
$g-kP<AP\le g$, while
$b'(1-\beta)\le g\le b'(1+\beta)$.
The precise standard inequalities used are
\[
 40\epsilon\le\delta<1/6,\quad
 0\le\beta\le\delta/2-18\epsilon,\quad
 E_0+Q\le\epsilon a,\quad
 kP\le kQ\le\epsilon b',\quad 0\le E_0,\quad Q_j\le Q.
 \tag{SW10.1}
\]
The distance $d$ is odd. The forbidden band about $a/b'$ and the
definition of the split index $\sigma$ imply
$p^d<(a/b')(1-2\delta)$ for $i<\sigma$, and
$p^d>(a/b')(1+2\delta)$ for $i\ge\sigma$.
The elementary inequalities
\[
 (1-2\delta)(1+\beta)(1+\delta)\le1-\epsilon,\qquad
 (1-\beta-\epsilon)(1-\delta)(1+2\delta)\ge1
\]
now give, respectively,
\[
 A(1+\delta)<A_0\quad(i<\sigma),\qquad
 A_0<A(1-\delta)\quad(i\ge\sigma).
 \tag{SW10.2}
\]
For the first, multiply by $Q_j$ and compare with
$a(1-\epsilon)<A_0Q_j$, retaining the strict floor inequality.
For the second use $AP>b'(1-\beta-\epsilon)$ and $A_0Q_j\le a$.
Finally $kQ\le\epsilon b'$ follows from $b'\ge k^3Q$,
$\epsilon\ge1/k$ and $k\ge1$.
Thus the exact coarse gap required by RL follows from literal
\texttt{Std}, including its excluded-band hypothesis.

\paragraph{SW11: scope and the next obligation.}
This completes the integral-form and summation mechanisms.
To obtain the full Type-I exceptional-function theorem one must still
supply the coarse rank/density input (including the lower-rank half),
the full numerical parameter package and the final comparison
$t_0\le t_I$. RL and SW12 below construct the actual cleaned lifts under
literal \texttt{Std}, including their raw-lift numerical conditions.
The common implication to RH and Dirichlet GRH
also retains the global obligations listed later.
